% META: {"id": "TSPE-DERIVATION-CONVEXITE-RE-C5", "chapitre": "TSPE-DERIVATION-CONVEXITE", "type_objet": "remediation", "capacites_codes": ["C5", "C6"], "status": "approved"}

%---------------------------------------------------------------
% FICHE DE REMEDIATION -- C5,C6 : Erreurs de lecture graphique et position courbe/tangente
%---------------------------------------------------------------

\subsection*{Rappel -- Erreurs de lecture graphique et position courbe/tangente}

\textbf{Erreur 1 :} Inverser la position courbe/tangente selon la convexite.

\textbf{Erreur 2 :} Croire que toute tangente est traversée (seuls les points d'inflexion traversent la tangente).

\textbf{Erreur 3 :} Confondre l'inegalite de convexite avec l'inegalite des accroissements.

\begin{exercice}{TSPE-DERCONV-RE-C5-EX1}{1}{12}
Soit $f(x) = \ln x$ sur $]0, +\infty[$.
\begin{enumerate}
  \item Determiner la tangente en $x = 1$.
  \item La courbe est-elle au-dessus ou au-dessous de cette tangente ?
  \item En quel point la courbe traverse-t-elle sa tangente ?
\end{enumerate}
\end{exercice}

% BEGIN-VERIFY
% from sympy import *
% x = symbols('x', positive=True)
% f = ln(x)
% fp = diff(f, x)
% assert fp.subs(x, 1) == 1
% assert f.subs(x, 1) == 0
% # tangente: y = x - 1
% g = f - (x - 1)
% gpp = diff(g, x, 2)
% assert gpp == -1/x**2  # < 0, concave => courbe sous la tangente
% assert g.subs(x, 2) == ln(2) - 1
% # f'' = -1/x^2 never zero, no inflexion point
% assert solve(diff(f, x, 2), x) == []
% END-VERIFY
